% ============================================================
%  THE ORTYX PROTOCOL
%  Part IV — Reconstruction
%  Chapter 16: The Reinterpretation Functor
% ============================================================

\chapter{The Reinterpretation Functor}
\label{ch:reinterpretation-functor}

\chapepigraph{Translation is not betrayal. It is the discovery
that the same structure was legible in two languages all along,
and that one of those languages makes the structure's
assumptions more visible than the other.}{Flyxion}

\noindent
The Interlude named what survived the decomposition of the
Phoenix corpus. Part~IV now performs the reconstruction:
the formal translation of survivable components into a
more stable semantic framework. The vehicle for that
translation is the reinterpretation functor:
\[
  \RSVPfunctor : \ThetaS \to \ThetaStar.
\]
This chapter specifies what the functor does and what
it does not do, before applying it to the survivable
components in Chapters~17--20.

\section{What the Functor Is}
\label{sec:functor-definition}

The reinterpretation functor is not a claim that the
Phoenix corpus is secretly a special case of RSVP, nor
that RSVP is more true than Phoenix. It is a structure-preserving
map from the survivable components of the Phoenix corpus
to a more stable expressive context: one in which the
components' assumptions are made explicit, their projection
dependencies are acknowledged, and their empirical contact
points are more clearly marked.

A functor, in the categorical sense, maps objects and
morphisms between categories while preserving composition
and identity. The reinterpretation functor $\RSVPfunctor$
maps:

\textit{Phoenix concepts} (the objects of $\ThetaS$)
to \textit{RSVP-compatible concepts} (the objects of
$\ThetaStar$), where the mapping preserves the structural
relationships between concepts --- the ways they imply,
constrain, and generate each other.

\textit{Phoenix arguments} (the morphisms of $\ThetaS$)
to \textit{RSVP-compatible arguments} (the morphisms
of $\ThetaStar$), where the mapping preserves the
inferential structure of arguments --- what they assume
and what they conclude.

The functor is not required to be injective (different
Phoenix concepts may map to the same RSVP concept) or
surjective (not all RSVP concepts need be images of
Phoenix concepts). It is required to be structure-preserving:
if two Phoenix concepts are related by a specific inferential
relationship in $\ThetaS$, their images must be related
by a corresponding relationship in $\ThetaStar$.

\section{What the Functor Is Not}
\label{sec:functor-not}

The functor is not a replacement of the Phoenix framework
by RSVP. The survivable components in $\ThetaS$ are
already grounded; the functor translates them, it does
not improve them. A jitter metric that is correctly
computed under the Marine algorithm remains correctly
computed after the functor maps it to an accessibility
entropy concept. The functor adds expressive resources
and makes assumptions explicit; it does not add empirical
grounding that was not there before.

The functor is not a proof that the translated components
are true in any stronger sense than they were before
translation. The translation makes the components more
auditable --- more clearly connected to their assumptions,
more clearly bounded in their scope --- but auditability
is not truth. The RSVP-translated components remain
subject to the same empirical requirements as the original
Phoenix components.

\begin{auditprop}
\label{ap:functor-scope}
The reinterpretation functor $\RSVPfunctor$ maps the
survivable components $\ThetaS$ of the Phoenix corpus
to a more stable expressive context $\ThetaStar$ by
making the components' projection dependencies explicit,
marking their empirical contact points more clearly,
and providing a vocabulary in which their boundary
conditions can be stated without the failure geometries
that afflicted the original formulation. The functor
does not add empirical grounding, improve predictive
accuracy, or establish the truth of any claim that was
not already established in $\ThetaS$.
\end{auditprop}

\section{The Projection Defect Reduction}
\label{sec:projection-defect-reduction}

The central virtue of the translated components in
$\ThetaStar$ over the original components in $\ThetaS$
is the reduction in projection defect. Recall from
Chapter~11 that the projection defect of a framework
is:
\[
  \ProjDef(\SysTriple)
  = \sup_{x \in X_{\mathrm{adm}}}
  \abs{\ConstrX(x) - \ConstrM(\pi(x))},
\]
measuring how much the constraints operative in the
full space fail to be reflected in the projected
representation.

The Phoenix corpus's projection defect is elevated
primarily because its claims about consciousness,
authenticity, and meaning are made at the level of
the full space (phenomenological and ontological claims
about the world) while the supporting evidence is at
the level of the projected representation (engineering
results about signal processing). The translation to
RSVP does not close this gap in the sense of providing
new evidence, but it closes it in a different sense:
the RSVP vocabulary makes the gap explicit rather than
obscuring it behind unified terminology.

In RSVP-compatible terms, a claim that would previously
read as: ``the Marine detects meaningful structure''
now reads as: ``the Marine detects low-entropy trajectory
regions in signal space, and whether low-entropy
trajectory regions correspond to meaningful structure
is an open empirical question that requires specification
of the meaning-accessibility relationship.'' The
second formulation is less rhetorically powerful but
more epistemically honest: it locates the projection
dependency explicitly rather than hiding it in
vocabulary continuity.

\section{The Translation Table}
\label{sec:translation-table}

The core translations that the functor performs are
summarized below. Each row maps a Phoenix concept or
claim to its RSVP-compatible counterpart.

\begin{center}
\begin{tabular}{p{0.42\textwidth}p{0.42\textwidth}}
\toprule
Phoenix $\ThetaS$ & RSVP $\ThetaStar$ \\
\midrule
Temporal coherence (jitter) & Trajectory constraint density:
low-jitter signals occupy low-$\AccEnt$ regions \\[6pt]
Marine salience filter & Constraint-governed gating: selection
of events consistent with admissible trajectories \\[6pt]
Harmonic conformity $\Gamma_n$ & Constraint propagation: harmonics
evolving within a shared admissible basin \\[6pt]
\MEMeight{} wave superposition & Field dynamics in the RSVP plenum:
wave interference as a parameterization of field-theoretic
memory \\[6pt]
Predictive processing alignment & Constraint-governed trajectory
selection: prediction error as trajectory revision signal \\[6pt]
Behavioral injection & Constraint field poisoning: manipulation
of the contextual field that governs admissible reasoning
trajectories \\[6pt]
Compression efficiency & Constraint discovery: compression that
exploits relational structure reduces redundancy by identifying
shared admissibility conditions \\
\bottomrule
\end{tabular}
\end{center}

\medskip

Each translation in the table preserves the core content
of the Phoenix concept while replacing vocabulary that
carries unacknowledged projection risks with vocabulary
that makes the projection dependencies explicit. The
RSVP column does not claim more than the Phoenix column;
it claims the same thing in terms that are more clearly
bounded.

\medskip

Chapters~17--20 apply this translation to the survivable
components in detail. Chapter~17 handles the signal
processing and salience results; Chapter~18 handles
the memory architecture; Chapter~19 handles the Nexus
framework as a new case study in constraint-preserving
computational memory; Chapter~20 synthesizes the
reconstructed components into a coherent alternative
paradigm statement.
